%% ma: L11-B11-0003
%% lop: 11
%% chuong: 4
%% bai: 11
%% dang: 11.11.4
%% loai: TL
%% muc: VD
%% dap_an: "-"
%% nguon: "de-cuong-giua-ky-1-toan-11-nam-2026-2027-truong-thpt-hoang-van-thu-ha-noi-2-8.pdf (M08, Đề cương giữa kì 1 Toán 11 THPT Hoàng Văn Thụ, Hà Nội), Phần 4 Bài 5"
%% kien_thuc: "Bài 10 (giao tuyến, giao điểm), Bài 11 (định lí về giao tuyến song song), định lí Thalès (lớp 8)"
%% trang_thai: chinh-thuc
%% ghi_chu: "Đã kiểm bằng tọa độ: K nằm trên đường thẳng AF, AK = 2AF (ngoài đoạn AF), KF/KA = 1/2. Dữ kiện G (trọng tâm tam giác ACD) không dùng đến trong các ý của đề. Đã duyệt 29/10/2026."
\begin{ex}
Cho hình chóp $S.ABCD$ có đáy $ABCD$ là hình bình hành. Gọi $E$ là điểm thuộc cạnh $CD$ sao cho $CD=3CE$, $F$ là điểm thuộc cạnh $SD$ sao cho $SF=\dfrac34SD$, $G$ là trọng tâm tam giác $ACD$.

a) Tìm giao tuyến của hai mặt phẳng $(SBE)$ và $(SAC)$.

b) Tìm giao tuyến của hai mặt phẳng $(SAD)$ và $(SBC)$; $(SAB)$ và $(SCD)$.

c) Tìm điểm $K$ là giao điểm của $AF$ và $(SBE)$. Tính tỉ số $\dfrac{KF}{KA}$.
\loigiai{a) Trong mặt phẳng đáy, gọi $H=BE\cap AC$. Khi đó $S$, $H$ là hai điểm chung, $(SBE)\cap(SAC)=SH$.

b) $(SAD)$, $(SBC)$ có điểm chung $S$ và lần lượt chứa $AD\parallel BC$, nên giao tuyến là đường thẳng qua $S$ song song với $AD$. Tương tự, $(SAB)\cap(SCD)$ là đường thẳng qua $S$ song song với $AB$.

c) Trong mặt phẳng đáy, gọi $P=BE\cap AD$. Khi đó $(SBE)\cap(SAD)=SP$, và $K=AF\cap SP$ (trong mặt phẳng $(SAD)$).

Vì $DE\parallel AB$ và $DE=\dfrac23AB$ nên $\dfrac{PD}{PA}=\dfrac{DE}{AB}=\dfrac23$, suy ra $PD=2AD$. Trong $\triangle SAD$, kẻ $FL\parallel AD$ ($L\in SP$): $\dfrac{FL}{PD}=\dfrac{SF}{SD}=\dfrac34$ nên $FL=\dfrac34\cdot2AD=\dfrac32AD$. Mà $FL\parallel AP$ với $AP=3AD$, nên $\dfrac{KF}{KA}=\dfrac{FL}{AP}=\dfrac{3/2}{3}=\dfrac12$ ($K$ nằm trên tia $AF$, ngoài đoạn $AF$).}
\end{ex}
